\section{연구 데이터 자산 구성}
\label{sec:research-assets}

\begin{itemize}
    \item \textbf{전체 구성 및 저장 위치}
    \begin{itemize}
        \item 본 연구용 데이터 자산은
        \texttt{RCDatasetCreation/dataset\_resources\_research}에 저장한다.
        \item 전체 용량은 약 4.3~GB이며, 배경 이미지와 투명 물체용
        메시로 구성한다.
        \item 메인 데이터 생성 설정은
        \texttt{configs/dataset\_prism\_main.yaml}에서 관리한다.
        \item 대용량 원본 자산은 Git 저장소에서 제외하고, 파일 해시와
        출처를 기록한 manifest만 버전 관리한다.
    \end{itemize}

    \item \textbf{배경 이미지}
    \begin{itemize}
        \item DIV2K 고해상도 이미지 총 900장을 사용한다.
        \item 초기 training pool은 800장이고, held-out test set은
        100장이다.
        \item 분할 목록은
        \texttt{background/train\_background.txt}와
        \texttt{background/test\_background.txt}에 기록한다.
        \item 각 이미지는 투명 물체 뒤의 유한 emissive plane에 texture로
        배치한다.
        \item plane의 깊이, 기울기, 밝기, crop, 회전, offset 및 좌우 반전을
        무작위화하여 배경 다양성을 확보한다.
        \item emissive plane을 사용함으로써 유효한 굴절 광선에 대한 정확한
        2차원 source coordinate를 계산하고, 물체 그림자나 배경 의존 조명의
        영향을 줄인다.
    \end{itemize}

    \item \textbf{투명 물체 메시}
    \begin{itemize}
        \item 정규화된 PLY 메시 총 150개를 사용한다.
        \item Objaverse에서 선별한 메시 40개와 직접 생성한 procedural 메시
        110개로 구성한다.
        \item procedural 메시에는 bottle, vase, bowl, lens/ellipsoid,
        prism 및 asymmetric ornament 계열의 solid/hollow 형상이 포함된다.
        \item 초기 training pool은 120개이고, held-out test set은 30개이다.
        \item 분할 목록은 \texttt{shape/train\_shape.txt}와
        \texttt{shape/test\_shape.txt}에 기록한다.
        \item PLY 파일에는 형상만 저장하며, 투명 재질은 메시 자체에 포함하지
        않는다.
        \item 렌더링 시 Mitsuba에서 dielectric 재질을 적용하고, 실험 설정에
        따라 IOR과 RGB transmission을 무작위로 선택한다.
        \item 동일한 형상에 다양한 광학 물성을 적용하면서 alpha,
        transmission, refractive correspondence 및 image formation 항의
        정확한 ground truth를 생성할 수 있다.
        \item Objaverse 원본 ID, 출처 URL과 라이선스는
        \texttt{shape/ATTRIBUTION.csv} 및 asset manifest에 기록한다.
    \end{itemize}

    \item \textbf{Train/Validation/Test 분할}
    \begin{itemize}
        \item training pool의 10\%를 고정 seed로 validation set에 분리한다.
        \item held-out test 목록은 validation 분할 과정에서 사용하지 않는다.
        \item 최종 메시 분할은 train 108개, validation 12개, test 30개이다.
        \item 최종 배경 분할은 train 720장, validation 80장, test 100장이다.
        \item 생성된 dataset manifest에 세 분할을 모두 기록한다.
        \item contract validator는 split 사이에 동일한 메시나 배경이 포함되면
        오류로 처리하여 data leakage를 방지한다.
    \end{itemize}

    \item \textbf{Paired-background 구성}
    \begin{itemize}
        \item 동일한 물체 형상, IOR, RGB transmission, 카메라 및 물체
        trajectory를 유지한 채 서로 다른 배경으로 반복 렌더링한다.
        \item 물리 조건을 고정하고 배경만 변경함으로써 모델의 배경 불변성과
        inverse refractive evidence의 유효성을 통제된 조건에서 평가한다.
        \item paired group 내부의 물체 pose, correspondence와 material
        parameter 일치 여부도 contract validator로 검사한다.
    \end{itemize}

    \item \textbf{Smoke-test용 소형 자산}
    \begin{itemize}
        \item 빠른 실행 검증용 자산은
        \texttt{RCDatasetCreation/dataset\_resources}에 별도로 저장한다.
        \item sphere와 octahedron 메시, 소형 planar background 4장 및 간단한
        3D tree scene을 포함한다.
        \item CPU smoke test, renderer 연동, 데이터 계약 및 split leakage
        검증에만 사용하며 본 학습 데이터에는 포함하지 않는다.
        \item smoke train/test fixture도 서로 겹치지 않도록 분리한다.
    \end{itemize}

    \item \textbf{재현성과 GitHub 배포 정책}
    \begin{itemize}
        \item 실제 연구 자산과 생성된 렌더 결과는 용량 및 제3자 배포 조건을
        고려하여 GitHub에 포함하지 않는다.
        \item 대신
        \texttt{asset\_manifests/prism\_research\_assets\_v1.json}에 각
        자산의 파일명, source, split, byte size, SHA-256 checksum 및
        license 정보를 기록한다.
        \item asset 준비 스크립트로 동일한 디렉터리와 split 목록을 재구성할
        수 있다.
        \item manifest freezing 및 contract validation 도구를 이용해 자산의
        동일성과 생성 데이터의 무결성을 검증한다.
    \end{itemize}
\end{itemize}
